% =====================================================================
%  ERRATA (opcional) -- Manual, seção 2.1.3
%  Só use se, depois de impresso/entregue, você encontrar erros.
%  Para ativar, retire o % da linha % =====================================================================
%  ERRATA (opcional) -- Manual, seção 2.1.3
%  Só use se, depois de impresso/entregue, você encontrar erros.
%  Para ativar, retire o % da linha % =====================================================================
%  ERRATA (opcional) -- Manual, seção 2.1.3
%  Só use se, depois de impresso/entregue, você encontrar erros.
%  Para ativar, retire o % da linha % =====================================================================
%  ERRATA (opcional) -- Manual, seção 2.1.3
%  Só use se, depois de impresso/entregue, você encontrar erros.
%  Para ativar, retire o % da linha \input{pretextuais/errata} no main.tex.
% =====================================================================
\begin{errata}
  \SingleSpacing
  \noindent\MakeTextUppercase{\imprimirautor}. \textbf{\imprimirtitulo}.
  \imprimirdata. Trabalho de Conclusão de Curso (\imprimirgrau) --
  \imprimirinstituicao, \imprimirlocal, \imprimirdata.

  \vspace{\baselineskip}
  \begin{table}[htb]
    \centering
    \begin{tabular}{cccc}
      \toprule
      \textbf{Folha} & \textbf{Linha} & \textbf{Onde se lê} & \textbf{Leia-se} \\
      \midrule
      10 & 2  & degau   & degrau   \\
      34 & 23 & atólogo & autólogo \\
      \bottomrule
    \end{tabular}
  \end{table}
\end{errata}
 no main.tex.
% =====================================================================
\begin{errata}
  \SingleSpacing
  \noindent\MakeTextUppercase{\imprimirautor}. \textbf{\imprimirtitulo}.
  \imprimirdata. Trabalho de Conclusão de Curso (\imprimirgrau) --
  \imprimirinstituicao, \imprimirlocal, \imprimirdata.

  \vspace{\baselineskip}
  \begin{table}[htb]
    \centering
    \begin{tabular}{cccc}
      \toprule
      \textbf{Folha} & \textbf{Linha} & \textbf{Onde se lê} & \textbf{Leia-se} \\
      \midrule
      10 & 2  & degau   & degrau   \\
      34 & 23 & atólogo & autólogo \\
      \bottomrule
    \end{tabular}
  \end{table}
\end{errata}
 no main.tex.
% =====================================================================
\begin{errata}
  \SingleSpacing
  \noindent\MakeTextUppercase{\imprimirautor}. \textbf{\imprimirtitulo}.
  \imprimirdata. Trabalho de Conclusão de Curso (\imprimirgrau) --
  \imprimirinstituicao, \imprimirlocal, \imprimirdata.

  \vspace{\baselineskip}
  \begin{table}[htb]
    \centering
    \begin{tabular}{cccc}
      \toprule
      \textbf{Folha} & \textbf{Linha} & \textbf{Onde se lê} & \textbf{Leia-se} \\
      \midrule
      10 & 2  & degau   & degrau   \\
      34 & 23 & atólogo & autólogo \\
      \bottomrule
    \end{tabular}
  \end{table}
\end{errata}
 no main.tex.
% =====================================================================
\begin{errata}
  \SingleSpacing
  \noindent\MakeTextUppercase{\imprimirautor}. \textbf{\imprimirtitulo}.
  \imprimirdata. Trabalho de Conclusão de Curso (\imprimirgrau) --
  \imprimirinstituicao, \imprimirlocal, \imprimirdata.

  \vspace{\baselineskip}
  \begin{table}[htb]
    \centering
    \begin{tabular}{cccc}
      \toprule
      \textbf{Folha} & \textbf{Linha} & \textbf{Onde se lê} & \textbf{Leia-se} \\
      \midrule
      10 & 2  & degau   & degrau   \\
      34 & 23 & atólogo & autólogo \\
      \bottomrule
    \end{tabular}
  \end{table}
\end{errata}
